% ============================================================
%  CV de Luís Otávio Silva Santos / Template ATS-friendly
%  Regras seguidas: coluna única, sem tabelas/listas de layout
%  crítico, contato no corpo, fonte Arial, bullets simples,
%  datas em formato consistente, texto selecionável.
%  Compilar com: tectonic -X compile <arquivo>.tex
% ============================================================
\documentclass[10pt,a4paper]{article}

\usepackage{fontspec}
% HyphenChar=None desativa a hifenização (XeTeX). Palavras técnicas
% nunca são quebradas no meio (melhor para ATS e para leitura).
\setmainfont{Arial}[HyphenChar=None]

\usepackage[a4paper,left=1.6cm,right=1.6cm,top=1.4cm,bottom=1.4cm]{geometry}
\usepackage{enumitem}
\usepackage{hyperref}
\usepackage{xcolor}
\usepackage[none]{hyphenat}

% Hifenização desativada: palavras nunca são quebradas no meio
% (importante para parsing de ATS e para legibilidade do texto).
\hyphenpenalty=10000
\exhyphenpenalty=10000
\sloppy

\definecolor{linknavy}{HTML}{14406B}

\hypersetup{
  colorlinks=true,
  urlcolor=linknavy,
  linkcolor=linknavy,
  pdfborder={0 0 0},
  pdfcreator={Tectonic},
}

\pagestyle{empty}
\setlength{\parindent}{0pt}
\setlength{\parskip}{4pt}
\setlength{\emergencystretch}{2em}

% ---- Cabeçalho de contato (no corpo, não em header/footer) ----
\newcommand{\cvheader}[3]{%
  % #1 título profissional  #2 linha de contato  #3 linha de links
  \begin{center}%
    {\LARGE\bfseries Luís Otávio Silva Santos}\\[2pt]
    {\normalsize #1}\\[3pt]
    {\small #2}\\[1pt]
    {\small #3}%
  \end{center}%
  \vspace{2pt}%
}

% ---- Seção (heading padrão, texto real) ----
\newcommand{\cvsection}[1]{%
  \begingroup
  \setlength{\parskip}{0pt}%
  \vspace{6pt}%
  {\normalsize\bfseries\MakeUppercase{#1}}\par
  \vspace{1.5pt}%
  \noindent\rule{\textwidth}{0.8pt}\par
  \endgroup
  \vspace{2pt}%
}

% ---- Experiência / formação: empregador, local, cargo, período ----
\newcommand{\cventry}[4]{%
  \vspace{3pt}%
  {\bfseries #1} \hfill #2\\*[1pt]
  {#3} \hfill {#4}\par
}

% ---- Projeto: nome + stack/stack-line + link opcional ----
\newcommand{\projentry}[3]{%
  % #1 nome   #2 stack   #3 link (texto exibido; vazio omite)
  \vspace{3pt}%
  {\bfseries #1}: {\itshape #2}%
  \ifx\relax#3\relax\else\hfill {\small\href{https://#3}{#3}}\fi
}

% ---- Bullets padronizados ----
\newenvironment{cvbullets}{%
  \begin{itemize}[leftmargin=1.15em,topsep=1.5pt,itemsep=1.1pt,parsep=0pt,partopsep=0pt,label=\textbullet]%
}{\end{itemize}}

% ---- Linha de competência (Skill: itens) ----
\newcommand{\skillline}[2]{%
  {\bfseries #1:} #2\par
}
